\chapter{Descent maps (Lemmas 10.4, 10.5, Prop 10.8)}
\label{chap:descent}

Descent maps remove the left-most column of $\mathcal{P}$, of $\mathcal{Q}$, or
of both, producing a PBP on a shifted shape. They are the central structural
tool of [BMSZb] Section~10.

\section{Single naive descent paints}

There are four single descent directions:
\begin{itemize}
\item D $\to$ C: cut col~0 of $\mathcal{Q}$ and shift $\mathcal{P}$ left by one column (refill top with $\bullet$).
\item C $\to$ D: cut col~0 of $\mathcal{P}$ (refill top/middle with $\bullet, s$).
\item B $\to$ M: cut col~0 of $\mathcal{Q}$ (refill $\mathcal{P}$, translate).
\item M $\to$ B: cut col~0 of $\mathcal{P}$.
\end{itemize}

\begin{definition}[D $\to$ C paints]
  \label{def:descentPaint_DC}
  \lean{descentPaintL_DC, descentPaintR_DC}
  \leanok
  \uses{def:PBP, def:dotScolLen}
  Let $c_L(j) := \mathrm{dotScolLen}(\tau.\mathcal{P}, j + 1)$ and $c_R(j) := \tau.\mathcal{Q}.\mathrm{colLen}(j)$.
  \begin{align*}
    (\nabla_{DC}\tau).\mathcal{P}(i, j) &= \begin{cases} \bullet & i < c_L(j) \\ \tau.\mathcal{P}(i, j + 1) & \text{otherwise} \end{cases} \\
    (\nabla_{DC}\tau).\mathcal{Q}(i, j) &= \begin{cases} \bullet & i < c_L(j) \\ s & c_L(j) \leq i < c_R(j) \\ \bullet & i \geq c_R(j) \end{cases}
  \end{align*}
\end{definition}

\begin{definition}[C $\to$ D paints]
  \label{def:descentPaint_CD}
  \lean{descentPaintL_CD, descentPaintR_CD}
  \leanok
  \uses{def:PBP, def:dotScolLen}
  $c_L(j) := \mathrm{dotScolLen}(\tau.\mathcal{P}, j)$, $c_R(j) := \tau.\mathcal{Q}.\mathrm{colLen}(j + 1)$.
  \[
    (\nabla_{CD}\tau).\mathcal{P}(i, j) = \begin{cases} \bullet & i < c_R(j) \\ s & c_R(j) \leq i < c_L(j) \\ \tau.\mathcal{P}(i, j) & i \geq c_L(j) \end{cases}
  \]
  $(\nabla_{CD}\tau).\mathcal{Q}$ is all $\bullet$ (D type has $\mathcal{Q}$-allowed $= \{\bullet\}$).
\end{definition}

\begin{definition}[B $\to$ M paints]
  \label{def:descentPaint_BM}
  \lean{descentPaintL_BM, descentPaintR_BM}
  \leanok
  \uses{def:PBP, def:dotScolLen}
  $c_L(j) := \mathrm{dotScolLen}(\tau.\mathcal{P}, j)$, $c_R(j) := \mathrm{dotScolLen}(\tau.\mathcal{Q}, j + 1)$.
  \begin{align*}
    (\nabla_{BM}\tau).\mathcal{P}(i, j) &= \begin{cases} \bullet & i < c_R(j) \\ s & c_R(j) \leq i < c_L(j) \\ \tau.\mathcal{P}(i, j) & i \geq c_L(j) \end{cases} \\
    (\nabla_{BM}\tau).\mathcal{Q}(i, j) &= \begin{cases} \bullet & i < c_R(j) \\ \tau.\mathcal{Q}(i, j + 1) & i \geq c_R(j) \end{cases}
  \end{align*}
\end{definition}

\begin{definition}[M $\to$ B paints]
  \label{def:descentPaint_MB}
  \lean{descentPaintL_MB, descentPaintR_MB}
  \leanok
  \uses{def:PBP, def:dotScolLen}
  $c_L(j) := \mathrm{dotScolLen}(\tau.\mathcal{P}, j + 1)$, $c_R(j) := \mathrm{dotScolLen}(\tau.\mathcal{Q}, j)$.
  \begin{align*}
    (\nabla_{MB}\tau).\mathcal{P}(i, j) &= \begin{cases} \bullet & i < c_L(j) \\ \tau.\mathcal{P}(i, j+1) & i \geq c_L(j) \end{cases} \\
    (\nabla_{MB}\tau).\mathcal{Q}(i, j) &= \begin{cases} \bullet & i < c_L(j) \\ s & c_L(j) \leq i < c_R(j) \\ \tau.\mathcal{Q}(i, j) & i \geq c_R(j) \end{cases}
  \end{align*}
\end{definition}

\section{Double descent $\nabla^2$}

\begin{definition}[Double descent D type]
  \label{def:doubleDescent_D}
  \lean{doubleDescent_D_paintL}
  \leanok
  \uses{def:descentPaint_DC, def:descentPaint_CD, def:dotScolLen}
  Composite D $\xrightarrow{\nabla_{DC}}$ C $\xrightarrow{\nabla_{CD}}$ D.
  After simplification:
  \[
    (\nabla^2_D \tau).\mathcal{P}(i, j) = \begin{cases}
      \bullet & i < \tau.\mathcal{Q}.\mathrm{colLen}(j + 1) \\
      s & \tau.\mathcal{Q}.\mathrm{colLen}(j + 1) \leq i < \mathrm{dotScolLen}(\tau.\mathcal{P}, j + 1) \\
      \tau.\mathcal{P}(i, j + 1) & i \geq \mathrm{dotScolLen}(\tau.\mathcal{P}, j + 1)
    \end{cases}
  \]
  (The resulting $\mathcal{Q}$ is all $\bullet$.)
\end{definition}

\begin{definition}[Double descent B type]
  \label{def:doubleDescent_B}
  \lean{doubleDescent_B_paintL, doubleDescent_B_paintR}
  \leanok
  \uses{def:descentPaint_BM, def:descentPaint_MB, def:dotScolLen}
  Composite B $\xrightarrow{\nabla_{BM}}$ M $\xrightarrow{\nabla_{MB}}$ B.
  \begin{align*}
    (\nabla^2_B \tau).\mathcal{P}(i, j) &= \begin{cases} \bullet & i < \mathrm{dotScolLen}(\tau.\mathcal{P}, j + 1) \\ \tau.\mathcal{P}(i, j + 1) & i \geq \ldots \end{cases} \\
    (\nabla^2_B \tau).\mathcal{Q}(i, j) &= \begin{cases} \bullet & i < \mathrm{dotScolLen}(\tau.\mathcal{P}, j + 1) \\ s & \mathrm{dotScolLen}(\tau.\mathcal{P}, j + 1) \leq i < \mathrm{dotScolLen}(\tau.\mathcal{Q}, j + 1) \\ \tau.\mathcal{Q}(i, j + 1) & i \geq \ldots \end{cases}
  \end{align*}
\end{definition}

\section{Lemmas 10.4, 10.5 — descent injectivity}

\begin{lemma}[Lemma 10.4 C — C$\to$D descent is injective]
  \label{lem:10_4_C}
  \lean{descentCD_injective, lemma_10_4_C}
  \leanok
  \uses{def:descentPaint_CD}
  Let $\mu_P, \mu_Q$ be shapes with $\mathrm{shiftLeft}(\mu_Q) \leq \mu_P$.
  The map $\nabla_{CD} : \PBP_C(\mu_P, \mu_Q) \to \PBP_D(\mu_P, \mathrm{shiftLeft}(\mu_Q))$ is injective.
\end{lemma}

\begin{lemma}[Lemma 10.4 B — B$\to$M + signature is injective]
  \label{lem:10_4_B}
  \lean{lemma_10_4_B, PBP.descent_inj_B}
  \leanok
  \uses{def:descentPaint_BM, def:signature_pbp, def:epsilon_pbp, def:dotScolLen}
  For $\gamma \in \{B^+, B^-\}$ and $\tau_1, \tau_2 \in \PBP_\gamma(\mu_P, \mu_Q)$
  with the interlacing shape hypothesis $\mathrm{dotScolLen}(\tau_i.\mathcal{P}, j) \geq \mathrm{dotScolLen}(\tau_i.\mathcal{Q}, j + 1)$:
  if $\nabla_{BM}\tau_1 = \nabla_{BM}\tau_2$, $\Sign(\tau_1) = \Sign(\tau_2)$, $\eps_{\tau_1} = \eps_{\tau_2}$,
  then $\tau_1 = \tau_2$.
\end{lemma}

\begin{lemma}[Lemma 10.5 M — M$\to$B descent is injective]
  \label{lem:10_5_M}
  \lean{descentMB_PBP, descentMB_injective, lemma_10_5_M}
  \leanok
  \uses{def:descentPaint_MB}
  Equal M$\to$B descent paints $\Rightarrow$ equal M-type PBPs (same shapes).
\end{lemma}

\begin{lemma}[Lemma 10.5 D — double descent + signature injective]
  \label{lem:10_5_D}
  \lean{doubleDescent_D_injective_on_PBPSet, lemma_10_5_D}
  \leanok
  \uses{def:doubleDescent_D, def:signature_pbp, def:epsilon_pbp}
  On $\PBP_D(\mu_P, \mu_Q)$, the map $\tau \mapsto (\nabla^2\tau, \Sign(\tau), \eps_\tau)$ is injective.
\end{lemma}

\section{Proposition 10.8 — descent cardinality relations}

We now turn these injectivity statements into cardinality equalities.

\begin{proposition}[Prop 10.8(a) C primitive]
  \label{prop:10_8_C_prim}
  \lean{card_C_eq_card_D_primitive, prop_10_8_C_primitive}
  \leanok
  \uses{lem:10_4_C}
  If $\mu_Q.\mathrm{colLen}(0) \geq \mu_P.\mathrm{colLen}(0)$ (\emph{primitive case}), then
  \[ |\PBP_C(\mu_P, \mu_Q)| = |\PBP_D(\mu_P, \mathrm{shiftLeft}(\mu_Q))|. \]
\end{proposition}

\begin{proposition}[Prop 10.8(b) C balanced]
  \label{prop:10_8_C_bal}
  \lean{card_C_eq_DD_plus_RC_balanced, prop_10_8_C_balanced}
  \leanok
  \uses{lem:10_4_C}
  If $\mu_P.\mathrm{colLen}(0) = \mu_Q.\mathrm{colLen}(0) + 1$ (\emph{balanced case}), the C$\to$D descent image is $\{\sigma \in \PBP_D \mid \mathrm{tailClass}(\sigma) \in \{\mathrm{DD}, \mathrm{RC}\}\}$:
  \[ |\PBP_C(\mu_P, \mu_Q)| = \#\{\mathrm{DD}\} + \#\{\mathrm{RC}\}. \]
\end{proposition}

\begin{proposition}[Prop 10.8(a) M primitive]
  \label{prop:10_8_M_prim}
  \lean{descent_bijective_primitive, prop_10_8_M_primitive}
  \leanok
  \uses{lem:10_5_M}
  If $\mu_P.\mathrm{colLen}(0) > \mu_Q.\mathrm{colLen}(0)$ (\emph{primitive}):
  \[ |\PBP_M(\mu_P, \mu_Q)| = |\PBP_{B^+}(\mathrm{shiftLeft}(\mu_P), \mu_Q)| + |\PBP_{B^-}(\mathrm{shiftLeft}(\mu_P), \mu_Q)|. \]
\end{proposition}

\begin{proposition}[Prop 10.8(b) M balanced]
  \label{prop:10_8_M_bal}
  \lean{descent_image_balanced, prop_10_8_M_balanced}
  \leanok
  \uses{lem:10_5_M}
  If $\mu_P.\mathrm{colLen}(0) = \mu_Q.\mathrm{colLen}(0) > 0$:
  \[ |\PBP_M(\mu_P, \mu_Q)| = |\PBP_{B^+}(\cdots)| + |\{\sigma \in \PBP_{B^-}(\cdots) : \mathrm{layerOrd}(\sigma.\mathcal{Q}(c_2 - 1, 0)) > 1\}|. \]
\end{proposition}

\chapter{Double-descent injectivity (Prop 10.9, Cor 10.10)}
\label{chap:ddescent}

\begin{proposition}[Prop 10.9 D — double descent is injective modulo $(\Sign, \eps)$]
  \label{prop:10_9_D}
  \lean{PBP.ddescent_inj_D, prop_10_9_D}
  \leanok
  \uses{lem:10_5_D, def:signature_pbp, def:epsilon_pbp}
  For D-type PBPs $\tau_1, \tau_2$ with identical shapes, identical $\nabla^2\tau$, identical $\Sign$, and identical $\eps$: $\tau_1 = \tau_2$.
\end{proposition}

\begin{proposition}[Prop 10.9 B — double descent injective]
  \label{prop:10_9_B}
  \lean{PBP.ddescent_inj_B, prop_10_9_B}
  \leanok
  \uses{lem:10_4_B, def:signature_pbp, def:epsilon_pbp}
  Analogue for $\gamma \in \{B^+, B^-\}$, under the interlacing constraint.
\end{proposition}

\begin{corollary}[Cor 10.10 D/B — PBPSet injectivity]
  \label{cor:10_10}
  \lean{cor_10_10_D, cor_10_10_B}
  \leanok
  \uses{prop:10_9_D, prop:10_9_B}
  The map $\tau \mapsto (\nabla^2\tau, \Sign(\tau), \eps_\tau)$ is injective on $\PBP_\gamma(\mu_P, \mu_Q)$ for $\gamma \in \{D, B^+, B^-\}$.
\end{corollary}
